\chapter{总结与展望}

\section{本文工作总结}

本文针对端侧AI应用开发中面临的模型部署效率低、推理性能不足等问题，以鸿蒙操作系统和MindSpore深度学习框架为技术平台，设计并实现了一个智能图像分类应用。本文的主要工作和创新点总结如下：

（1）深入分析了鸿蒙系统的分布式架构和ArkTS开发框架，明确了端侧AI应用在鸿蒙系统上的开发技术路线。通过研究鸿蒙的Ability模型、NAPI接口机制和ArkUI声明式开发范式，提出了"训练端优化+端侧推理+声明式UI"的三层技术架构，为鸿蒙AI应用开发提供了系统性的技术方案。

（2）基于MindSpore框架完成了ResNet-50图像分类模型的训练与优化。在CIFAR-10数据集上，原始FP32模型达到了92.3\%的分类准确率。针对端侧设备计算资源有限的特点，采用INT8后训练量化和知识蒸馏两种轻量化策略，将模型体积从97.8MB压缩至24.5MB，推理延迟从112ms降低至45ms，精度仅下降1.2\%，实现了精度与效率的良好平衡。

（3）设计并实现了鸿蒙端侧推理引擎。通过NAPI接口封装MindSpore Lite C++推理库，设计了Model类的加载、预测和释放接口，实现了图像预处理（缩放、归一化、格式转换）和模型推理的完整流水线。推理引擎支持FP32和INT8两种推理模式，可根据设备性能自动选择最优模式。

（4）基于ArkTS框架设计了直观友好的用户交互界面。应用包含主页面、分类结果页面和历史记录页面三个核心页面，支持图片选取、实时分类、置信度可视化和历史记录管理等功能。界面采用声明式开发范式，代码结构清晰，可维护性强。

（5）在华为Mate 60 Pro真机上进行了全面的系统测试。功能测试表明所有核心功能正常工作；性能测试显示单张图片推理耗时约45ms，内存占用约120MB，满足实时性要求；用户体验测试SUS评分为82.5分，表明应用具有良好的可用性。

\section{存在的不足}

尽管本文在基于鸿蒙系统与MindSpore的端侧AI应用开发方面取得了一定的成果，但仍存在以下不足：

（1）分类类别有限。本文仅使用CIFAR-10数据集进行训练和测试，仅支持10个类别的图像分类，实际应用场景中的类别数量远超于此。扩展到ImageNet的1000个类别或自定义类别时，模型精度和推理速度可能面临更大的挑战。

（2）模型结构单一。本文仅使用了ResNet-50一种网络结构，未对其他轻量化网络（如EfficientNet、GhostNet等）在鸿蒙端侧的推理性能进行对比分析。

（3）分布式能力利用不足。鸿蒙系统的核心优势之一是分布式能力，本文的应用仅在单设备上运行，未充分利用鸿蒙的分布式软总线、分布式数据管理等能力实现跨设备AI协同。

（4）动态更新能力缺失。当前应用的模型文件打包在应用内部，无法在不更新应用的情况下更新模型。在实际应用中，模型可能需要定期更新以适应新的分类需求。

\section{未来工作展望}

针对上述不足，未来的研究工作可以从以下几个方面展开：

（1）扩展分类类别和模型能力。将分类类别从10类扩展到100类甚至1000类，引入目标检测、语义分割等更复杂的视觉任务，提升应用的实用性。同时，探索增量学习方法，使模型能够在端侧持续学习新的类别。

（2）探索更多轻量化模型。对比研究EfficientNet、GhostNet、ShuffleNetV2等轻量化网络在鸿蒙端侧的推理性能，寻找精度和效率的最优平衡点。研究神经架构搜索（NAS）技术在端侧模型设计中的应用。

（3）利用鸿蒙分布式能力。研究基于分布式软总线的跨设备模型推理方案，将计算密集型的推理任务卸载到附近的高性能设备上执行，实现设备间的AI能力共享。利用分布式数据管理实现模型和推理结果的跨设备同步。

（4）实现模型动态更新。设计基于鸿蒙应用市场的模型热更新机制，支持在不更新应用的情况下下载和替换模型文件，实现模型的持续迭代优化。

（5）增强隐私保护能力。研究联邦学习技术在端侧AI中的应用，在保护用户数据隐私的前提下实现模型的协同训练和优化，为端侧AI的隐私保护提供技术保障。

总之，随着鸿蒙系统生态的不断成熟和端侧AI技术的持续发展，基于鸿蒙系统的端侧AI应用将拥有广阔的应用前景。本文的研究工作为该领域的探索提供了有益的参考，期望未来能够看到更多高质量的端侧AI应用在鸿蒙生态中涌现。